\begin{textsample}{NEG Dim 2 – 2011\_c6 – Score -5.00 – 2011\_c6/enade\_2011\_eng\_g\_8\_80de7597d1.txt}  \label{ex:f2_neg_026}
CAPÍTULO 3 ANÁLISE TÉCNICA DA PROVA Este capítulo tem por objetivo apresentar o desempenho dos estudantes concluintes de Engenharia – Grupo VIII no ENADE/2011.
Para isso, foram calculadas as estatísticas básicas da prova em seu todo, bem como as estatísticas dos componentes relacionadas à Formação Geral, ao de Conhecimento Específico da Área e das questões discursivas isoladamente.
Nas tabelas, são apresentados o tamanho da população inscrita e de presentes, e as seguintes estatísticas das notas⁹ : média do desempenho na prova, erro padrão da média, desvio padrão, nota mínima, mediana e nota máxima.
As estatísticas apresentadas neste capítulo contemplam o total de estudantes concluintes da área de Engenharia – Grupo VIII em 2011 do Brasil e, separadamente, por Grande Região.
Foram calculadas tendo -se em vista as seguintes agregações : ( a ) as Grandes Regiões e o país como um todo ; ( b ) a Categoria Administrativa ; e ( c ) a Organização Acadêmica.
Em relação aos gráficos de distribuição de notas, o intervalo considerado foi de dez unidades, aberto à esquerda e fechado à direita, com exceção do primeiro intervalo, [ 0 ; 10 ], fechado em ambos os extremos.
Para os gráficos de distribuição das notas das questões discursivas, foram consideradas mais duas categorias : questão em branco e nota zero. 3.1 ESTATÍSTICAS BÁSICAS DA PROVA 3.1.1 Estatísticas Básicas Gerais A Tabela 3.1 apresenta as estatísticas básicas da prova por grande Região.
A população total de inscritos foi de 2.269.
Destes, 1.920 estiveram presentes, sendo 15,4 \% o índice de não comparecimento.
A Região de maior abstenção foi a Norte ( 30,0 \% ) e a de menor abstenção foi a Centro-Oeste ( 8,0 \% ). ⁹ Essas estatísticas e outras estão definidas no Capítulo 1.
A média das notas da prova como um todo ( nas seções seguintes serão analisados os Componentes de Formação Geral e de Conhecimento Específico ) foi 41,2, sendo que os alunos da região Nordeste obtiveram a média mais baixa ( 27,2 ) e os da região Sudeste obtiveram a média mais alta ( 50,6 ).
As demais médias foram : 39,6 na região Norte, 43,1 na região Sul e 46,9 na região Centro-Oeste.
O desvio padrão para o Brasil como um todo foi 18,0, sendo o maior desvio padrão encontrado na região Nordeste ( 20,0 ) e o menor na região Norte ( 13,0 ), indicando uma menor dispersão das notas desta última região.
A região que obteve a maior nota máxima foi a Sul ( 80,4 ), ao passo que a região que atingiu a menor nota máxima foi a Norte ( 74,0 ).
A mediana do Brasil como um todo foi 44,0, sendo a maior mediana obtida na região Sudeste ( 53,4 ) e a menor obtida na Nordeste ( 27,4 ).
A nota mínima foi zero em todas as regiões, sem exceção.
O comportamento das notas dos estudantes de todo o Brasil pode ser observado no Gráfico 3.1 que apresenta um histograma com a distribuição das mesmas.
Essa é uma distribuição unimodal com moda no intervalo ( 40 ; 50 ].
O coeficiente de assimetria da distribuição das notas é negativo ( -0,58 ).
As distribuições por Grande Região também apresentam assimetria negativa, concentração pouco maior do lado direito do histograma e mais espalhada do lado esquerdo.
A única exceção é a região Nordeste, com coeficiente de assimetria positivo muito próximo de zero ( 0,08 ).
Os Gráficos 3.2, 3.3 e 3.4 apresentam informações referentes à média da nota final dos Participantes, desagregando os resultados de acordo com, respectivamente, as Grandes Regiões do país, a Categoria Administrativa e a Organização Acadêmica.
Os gráficos apresentam o valor da média das notas como uma barra e os extremos do intervalo de confiança de 95 \% como linhas verticais unidas por uma linha horizontal na forma da letra H maiúscula.
Considerando -se o gráfico de notas segundo Grande Região ( Gráfico 3.2 ), observa -se que existe diferença estatisticamente significativa ao nível de 95 \% entre a menor média obtida na região Nordeste ( 27,2 ) e todas as demais regiões, em particular com a maior média, obtida na região Sudeste ( 50,6 ).
Levando -se em conta os estudantes por Categorias Administrativas ( Gráfico 3.3 ), observa -se que não existe diferença estatisticamente significativa entre as médias das notas das IES Públicas e Privadas, constatado pela diferença de 0,3 entre as médias dos alunos de IES Públicas e os alunos de IES Privadas.
Tendo como foco o Gráfico 3.4, que apresenta as notas médias das provas segundo Organização Acadêmica, contata -se que existe diferença estatisticamente significativa ao nível de 95 \% nas médias das notas dos estudantes provenientes de Centros Universitários em relação às demais Organizações Acadêmicas.
A maior média foi obtida pelos estudantes de Faculdades ( 42,5 ), e a menor, pelos de Centros Universitários ( 35,2 ).
Estudantes de Universidades obtiveram média 41,1. 3.1.2 Estatísticas Básicas no Componente de Formação Geral A Tabela 3.2 apresenta as estatísticas básicas em relação ao componente da prova que avalia a Formação Geral dos estudantes concluintes.
Os alunos de todo Brasil obtiveram desempenho médio de 51,5.
Quanto à variabilidade, o desvio padrão das notas dos estudantes do Brasil como um todo foi 21,3.
A maior média foi obtida na região Centro-Oeste ( 56,7 ), e a menor na região Nordeste ( 40,3 ).
As demais médias foram : 54,6 na região Norte, 56,1 na região Sudeste e 51,5 na região Sul.
Já o maior desvio padrão foi obtido na região Nordeste ( 27,1 ) e o menor na região Centro-Oeste ( 17,5 ).
Os demais desvios padrões foram : 17,6 na região Norte, 19,2 na região Sudeste e 18,9 na região Sul.
A maior nota no Componente de Formação Geral da prova do ENADE/2011 foi obtida por pelo menos um aluno das regiões Norte e Sudeste ( 98,0 ) enquanto que a menor nota máxima foi obtida na região Sul ( 85,5 ).
Nas outras regiões as notas máximas foram : 88,5 na região Nordeste e 94,0 na região Centro-Oeste.
A mediana do Brasil como um todo foi 55,5, sendo a menor mediana encontrada na região Nordeste ( 45,0 ) e a maior encontrada na região Sudeste ( 58,5 ).
A nota mínima nesta parte foi zero em todas as regiões, sem exceção.
O Gráfico 3.5 propicia a avaliação do desempenho dos estudantes no componente de Formação Geral a partir do histograma da distribuição das notas correspondentes.
A distribuição é unimodal, com moda em ( 50 ; 60 ], enquanto na prova como um todo a moda foi alcançada no intervalo ( 40 ; 50 ].
Nota -se, ainda, que no gráfico 3.5 as notas apresentam uma maior dispersão do que no Gráfico 3.1 ( distribuição das notas da prova ), confirmado pela comparação dos desvios padrões : 18,0 para a nota da prova como um todo e 21,3 para o Componente de Formação Geral.
Para o Componente de Formação Geral, o coeficiente de assimetria da distribuição das notas dos estudantes, como na prova como um todo, também é negativo ( -0,82 ).
A distribuição apresenta uma leve concentração à direita e cauda maior à esquerda, reforçada pelo máximo local no primeiro intervalo, [ 0 ; 10 ].
Em todas as Grandes Regiões os histogramas também possuem assimetria levemente negativa ( entre – 0,27 e – 1,04 ).
Nos Gráficos 3.6, 3.7 e 3.8 são apresentadas as informações referentes ao desempenho dos concluintes no Componente de Formação Geral, em diferentes agregações : Grande Região do país, Categoria Administrativa e Organização Acadêmica.
Observa -se pelo Gráfico 3.6 que existe diferença estatisticamente significativa entre as médias das notas no Componente de Formação Geral da região Nordeste ( 40,3 ) em relação às demais Grandes Regiões do país.
Vemos que o intervalo de confiança mais largo ( 5,3 ) é o da região Nordeste ; já o intervalo mais estreito ( 3,5 ) é observado nas regiões Sul e Norte.
Este fato está relacionado tanto com o tamanho da população envolvida quanto com o desvio padrão da média.
No Gráfico 3.7, que representa as notas médias no Componente de Formação Geral segundo Categoria Administrativa do país, observa -se que não existe diferença estatisticamente significativa entre as médias.
Os concluintes das IES Públicas obtiveram uma média menor ( 51,3 ) do que os das IES Privadas ( 52,9 ).
Considerando -se o tipo de Organização Acadêmica, nota -se, no Gráfico 3.8, não existe diferença estatisticamente significativa entre a maior e a menor média.
Nos Centros Universitários ( 55,4 ), a média é mais elevada do que em Faculdades ( 53,8 ) e Universidades ( 51,2 ). 3.1.3 Estatísticas Básicas do Componente de Conhecimento Específico A Tabela 3.3 apresenta as estatísticas básicas referentes ao Componente de Conhecimento Específico da área de Engenharia – Grupo VIII.
A média do desempenho dos alunos do Brasil como um todo foi 37,8.
A maior média foi obtida na região Sudeste ( 48,8 ), e a menor, na região Nordeste ( 22,9 ).
As demais médias foram : 34,5 na região Norte, 40,3 na região Sul e 43,6 na região Centro-Oeste.
Quanto à variabilidade das notas, o desvio padrão do Brasil como um todo foi 19,3, sendo o maior desvio padrão observado na região Nordeste ( 20,1 ) e o menor nas regiões Norte ( 14,5 ).
Os demais desvios foram : 16,6 da região Sul, 17,4 da região Sudeste e 16,2 da região Centro-Oeste.
A mediana das notas dos estudantes de todo o Brasil foi 40,5.
A maior mediana ocorreu na região Sudeste ( 51,0 ) e a menor na região Nordeste ( 20,2 ).
As demais medianas foram 35,4 na região Norte, 43,0 na região Sul e 44,5 na região Centro-Oeste.
A nota máxima do Brasil como um todo foi 83,1, sendo obtida por pelo menos um aluno da região Sudeste.
As demais notas máximas foram : 78,6 na região Norte, 76,6 na região Nordeste, 82,9 na região Sul e 79,3 na região Centro-Oeste.
A nota mínima foi zero em todas as regiões, sem exceção.
Assim como os Gráficos 3.1 e 3.5, o Gráfico 3.9, apresentado a seguir, proporciona uma avaliação do desempenho de concluintes em relação ao Componente de Conhecimento Específico com um histograma da distribuição das notas correspondentes.
Dentre as três distribuições apresentadas, esta é a mais concentrada nas notas baixas.
Esta também é uma distribuição unimodal, e o grupo modal é o ( 40 ; 50 ].
O coeficiente de assimetria da distribuição das notas do Componente de Conhecimento Específico é negativo ( – 0,35 ).
Nota -se pelo histograma ( Gráfico 3.9 ) que esta é uma distribuição aproximadamente simétrica, mas com cauda mais pesada do lado esquerdo, reforçada pelo máximo local no primeiro intervalo, [ 0 ; 10 ].
Já na região Nordeste, o coeficiente de assimetria é positivo, igual a 0,31, evidenciando que a distribuição das notas dos estudantes dessa região, no Componente de Conhecimento Específico, tem cauda mais pesada à direita.
Os Gráficos 3.10, 3.11 e 3.12 apresentam uma comparação dos resultados em relação à Grande Região do país, à Categoria Administrativa e à Organização Acadêmica, agora levando em conta o desempenho de alunos no Componente de Conhecimento Específico da prova.
Pelo Gráfico 3.10, observa -se que não existe diferença estatisticamente significativa entre as médias das notas no Componente de Conhecimento Específico apenas na comparação entre as regiões Sul e Centro-Oeste.
A região Sudeste foi a que obteve média mais elevada, enquanto a região Nordeste foi a região que obteve média mais baixa.
Quanto à Categoria Administrativa ( gráfico 3.11 ), observa -se um comportamento semelhante àquele da parte de Formação Geral e à prova como um todo, ou seja, não existe diferença estatisticamente significativa entre as médias das IES Públicas e Privadas.
Ao contrário, apesar de apresentar médias de valores próximos, a maior média foi obtida por alunos de IES Públicas de ensino.
Quanto ao Gráfico 3.12, observa -se, mais uma vez, que existe diferença estatisticamente significativa ao nível de 95 \% entre as notas no Componente de Conhecimento Específico dos Centros Universitários ( 28,5 ) em relação aos demais tipos de Organização Acadêmica.
Sendo que a média dos concluintes das Faculdades ( 38,7 ) foi maior do que de Universidades ( 37,7 ) e de Centros Universitários. 3.2 ANÁLISE DAS QUESTÕES OBJETIVAS 3.2.1 Componente de Formação Geral A Tabela 3.4 apresenta as estatísticas básicas relativas às oito questões objetivas do componente da prova que abrange a Formação Geral dos estudantes.
A média do Brasil foi 51,2.
A menor média foi encontrada na região Nordeste ( 42,8 ) e a maior na região Sudeste ( 55,3 ).
As demais médias foram : 54,1 na região Norte, 50,9 na região Sul e 54,0 na região Centro-Oeste.
O desvio padrão do Brasil foi 22,3, sendo o maior desvio padrão encontrado na região Nordeste ( 28,3 ) e o menor na região Centro-Oeste ( 19,0 ).
Os demais desvios foram : 19,9 na região Norte, 21,1 na região Sudeste e 19,3 na região Sul.
A mediana do Brasil foi 50,0.
A menor mediana foi encontrada a mesma do Brasil encontrada nas regiões Norte, Nordeste e Sul.
As regiões Sudeste e Centro-Oeste apresentaram a mesma mediana de 62,5.
As notas mínimas ( 0,0 ) foram iguais para todas as regiões.
A nota máxima do Brasil foi 100,0, encontrada em quase todas as regiões, a exceção foi a região Nordeste com 87,5.
A Tabela 3.5 apresenta os índices de facilidade e discriminação ( ponto bisserial ) das questões objetivas do Componente de Formação Geral.
Quanto ao índice de facilidade, foram usadas as seguintes cores para diferenciar o nível de dificuldade da questão : • Azul para as questões classificadas com índice muito fácil ( índice \&gt ; =0,86 ), verde para as questões classificadas com índice fácil ( 0,61 a 0,85 ), amarelo para as questões classificadas com médio ( 0,41 a 0,60 ), vermelho para as questões classificadas com difícil ( 0,16 a 0,40 ) e roxo para as questões classificadas com muito difícil ( \&lt ; =0,15 ).
Já quanto ao índice de discriminação, foram usadas as seguintes cores para qualificar a questão : • As questões classificadas com índice fraco receberam a cor vermelho ( índice \&lt ; =0,19 ), as classificadas com médio receberam a cor amarelo ( 0,20 a 0,29 ), as classificadas com bom receberam a cor verde ( 0,30 a 0,39 ) e as classificadas com muito bom ( \&gt ; =0,40 ) receberam a cor azul.
As questões objetivas do componente de Formação Geral, segundo o índice de facilidade, foram assim avaliadas : das oito questões, nenhuma teve o índice de facilidade classificado como muito fácil.
Quatro questões foram tidas como fáceis, por terem índice de acertos situado na faixa entre 0,61 e 0,85 ( de 61,0 \% a 85,0 \% de acertos ).
Duas questões foram consideradas de dificuldade média, situando -se no intervalo entre 0,41 e 0,60 do índice de facilidade, ou seja, houve entre 41,0 \% e 60,0 \% de acertos.
Outras duas questões foram classificadas como difícil, situando -se no intervalo entre 0,16 e 0,40.
Nenhuma questão apresentou menos de 15 \% de acertos, razão pela qual seria classificada como muito difícil.
Como já comentado, para análise das questões objetivas relativas à Formação Geral segundo o poder de discriminação, utilizou -se, o índice de discriminação ponto bisserial.
Nesta análise as questões foram assim avaliadas : cinco das oito questões apresentaram índices acima de 0,40 e, assim, foram classificadas com índice muito bom para esse grupo de alunos.
As outras três questões foram classificadas boas segundo o índice de discriminação, entre 0,30 e 0,39.
O índice de facilidade variou de 0,18 a 0,77, e o de discriminação, de 0,31 a 0,57.
As questões com índices de discriminação muito bom, de números 1, 2, 3, 5 e 6, figuraram entre as mais fáceis desse conjunto : quatro questões classificadas na categoria fácil ( questões 1, 3, 5 e 6 ) do índice de facilidade e uma questão, classificada na categoria médio ( questão 2 ).
Em particular, a questão 3 foi a que apresentou maior poder discriminatório, com índice 0,57, e dentre a mais fácil, com uma proporção de 0,73 de acertos.
A questão de número 8 apresentou índice de facilidade 0,18, ou seja, um quantitativo pequeno de estudantes conseguiu resolvê -la, dentro do universo de participantes.
No entanto, seu índice de discriminação foi bom.
As oito questões foram classificadas quanto o índice de discriminação como bom e muito bom.
O Gráfico 3.13, para exemplificar, analisa o comportamento da questão de número 3 de Formação Geral.
Trata -se de uma das questões mais fáceis e a que obteve o maior índice de discriminação dessa parte da prova.
Neste gráfico, cada uma das cinco curvas representa o percentual de respostas em determinada alternativa da questão, em função da nota dos estudantes nesta parte da prova ( Formação Geral/Múltipla Escolha ), antes de possíveis eliminações pelo critério do ponto bisserial.
A curva em vermelho corresponde à alternativa E, a correta para esta questão.
Assim, observa -se que entre os estudantes com um número menor de acertos, nessa parte do exame, a situação mais frequente foi a escolha de uma das alternativas incorretas : a alternativa A ( em azul ) ou D ( em roxo ).
Na medida em que a nota aumenta, indicando desempenho melhor nesta parte da prova, aumenta concomitantemente a proporção de estudantes que selecionaram a alternativa correta E, atingindo 100 \% para apenas os que acertaram todas as questões.
Essa análise permite verificar como a questão discriminou os grupos de desempenho, justificando o alto índice obtido na questão.
Os gráficos relativos às demais questões de Formação Geral constam do Anexo I. 3.2.2 Componente de Conhecimento Específico A Tabela 3.6 apresenta as estatísticas básicas em relação às questões objetivas do Componente de Conhecimento Específico da prova, por Grande Região.
A média do Brasil deste componente foi de 41,5.
A menor média foi observada na região Nordeste ( 25,7 ) e a maior na região Sudeste ( 53,2 ).
O desvio padrão de todo o Brasil foi 20,9, sendo o menor desvio padrão encontrado na região Norte ( 15,9 ) e o maior na região Nordeste ( 22,4 ).
A nota mínima de todo o Brasil foi 0,0, sem exceção.
A mediana do Brasil foi 42,9 e a região Nordeste apresentou a menor mediana ( 23,8 ).
Nas demais regiões, as medianas foram : 38,1 na região Norte, 57,1 na região Sudeste e 47,6 nas regiões Sul e Centro-Oeste.
A nota máxima do Brasil e das regiões Sudeste e Sul foi 90,5 e nas regiões Norte, Nordeste e Centro-Oeste foi 85,7.
A Tabela 3.7 apresenta os índices de facilidade e discriminação ( ponto bisserial ) das questões objetivas do Componente de Conhecimento Específico.
Para facilitar a diferenciação das questões usaram -se as mesmas cores da Tabela 3.5 para as diferentes classificações dos índices de facilidade e de discriminação.
Dentre as questões objetivas da parte da prova relativa ao Componente de Conhecimento Específico, não houve questão anulada pela Comissão.
Desse modo, a classificação quanto ao índice de facilidade foi estabelecida com base em 27 questões.
A partir dos índices obtidos, pode -se concluir que a maioria das questões objetivas da prova foi considerada de difícil a médio.
Das 27 questões, quatro foram classificadas como muito difíceis, abaixo de 0,15 do índice de facilidade, nove como difíceis, entre 0,16 e 0,40, e outras nove como médio, entre 0,41 e 0,60, e cinco como fácil, na faixa de 0,61 a 0,85 do índice de facilidade.
Nenhuma questão foi classificada como muito fácil.
Já quanto aos índices de discriminação das questões objetivas do Componente de Conhecimento Específico da prova, tem -se como resultado a seguinte classificação : quatro das 27 questões válidas foram consideradas como boas, enquanto nove delas tiveram índice de discriminação muito bom.
Assim, para quase metade das questões – 13 em 27 – os índices de discriminação foram bons ou muito bons.
Dentre as demais, oito delas foram classificadas como médias e outras seis como fracas, sendo quatorze, por conseguinte, a quantidade de questões nos dois patamares mais baixos de discriminação.
Constata -se, assim, que a prova – no que se refere ao componente de Conhecimento Específico – possuía boa capacidade de discriminar entre aqueles que dominam ou não o conteúdo.
Dentre as questões que alcançaram os maiores índices de discriminação, as de números 10, 18, 20, 22, 23, 28, 29, 33 e 34 classificadas com índice muito bom, situando -se no intervalo de 0,40 a 0,48 do índice, cinco delas ( questões 18, 22, 23, 33 e 34 ) foram classificadas na categoria fácil e quatro ( 10, 20, 28 e 29 ) como média, quanto ao índice de facilidade.
A questão de número 16 foi a mais difícil dentre as 27 questões específicas válidas, com baixo índice de facilidade, apenas 11,0 \% de acertos.
Essa questão apresentou poder discriminatório igualmente baixo, 0,13, o que comprova ter sido esta a mais difícil para os estudantes.
Destaca -se, também, a questão 30, com índice de facilidade 0,12, o que, em termos percentuais, corresponde a 12,0 \% de estudantes que responderam acertadamente, obtendo, ainda, 0,07 de índice de discriminação.
Tais questões foram, portanto, duas das mais difíceis da prova.
Pelo critério do ponto bisserial, as questões 16 e 30 foram eliminadas do cômputo da nota final.
Além destas duas, as demais questões com índice fraco de discriminação, questões 9, 14, 15 e 19 também não foram computadas.
A título de exemplo das análises do comportamento das questões objetivas, o Gráfico 3.14 analisa a questão 10 do Componente de Conhecimento Específico.
Esta foi uma questão classificada pelo índice de facilidade como média da prova, apresentando índice de facilidade 0,57, ou seja, 57,0 \% dos estudantes assinalaram acertadamente a opção A, correspondente ao gabarito.
Seu índice de discriminação foi igual a 0,48, classificado como muito bom.
Neste gráfico, cada uma das cinco curvas representa o percentual de respostas em determinada alternativa da questão 10, em função da nota dos estudantes nesta parte da prova, antes de possíveis eliminações de questões pelo critério do ponto bisserial.
A alternativa correta A, representada no gráfico pela curva em azul, foi escolhida em maiores proporções pelos alunos com desempenho melhor nesta parte da prova.
Já as alternativas incorretas, também denominadas distratores, foram selecionadas principalmente por aqueles com notas mais baixas.
Aqueles com nota zero, na sua quase totalidade deixaram esta questão em branco ou marcaram mais de uma alternativa, comportamento considerado inválido.
A proporção de alunos que selecionou a resposta correta A aumenta gradativamente, chegando a atingir 100 \% para 18 ou mais acertos, enquanto a proporção dos que escolheram alternativas incorretas decai, a partir das primeiras notas não nula, como função da nota nesta parte da prova.
Os gráficos relativos às demais questões do componente de Conhecimento Específico constam do Anexo I. 3.3 ANÁLISE DAS QUESTÕES DISCURSIVAS 3.3.1 Componente de Formação Geral As análises dos resultados de desempenho dos estudantes de Engenharia – Grupo VIII nas duas questões discursivas relativas à Formação Geral encontram -se na Tabela 3.8 e no Gráfico 3.15.
Na Tabela 3.8 observa -se que as notas médias foram ligeiramente mais altas nesse conjunto de questões do que no das objetivas.
Os estudantes de todo o Brasil obtiveram, em Formação Geral, média 51,2 nas questões objetivas e 51,9 nas questões discursivas.
A mediana confirma o desempenho um pouco melhor dos alunos de todo o Brasil nas questões discursivas do Componente de Formação Geral, enquanto essa estatística foi de 50,0 para as questões objetivas, para as questões discursivas a mediana foi 55,0.
Pode -se, também, notar um aumento do desvio padrão de 22,3 nas questões objetivas do Componente de Formação Geral dos alunos de todo o Brasil, para 29,9 nas questões discursivas do mesmo componente.
A maior mediana foi encontrada na região Centro-Oeste ( 65,0 ) e a menor na região Nordeste ( 40,0 ).
As demais medianas foram : 60,0 nas regiões Norte e Sudeste e 55,0 na região Sul.
A nota máxima ( 100,0 ) e a mínima ( 0,0 ) foram obtidas por alunos de todas as regiões do Brasil, sem exceção.
O Gráfico 3.15 representa a distribuição das notas nas questões discursivas no Componente de Formação Geral.
A moda desta distribuição ocorre no primeiro intervalo, ( 0 ; 10 ], devido à grande quantidade de notas zero.
Desconsiderando esse intervalo, a moda ocorre em ( 70 ; 80 ], com mais alunos tendo notas abaixo deste intervalo ( assimetria à esquerda ).
O coeficiente de assimetria da distribuição das notas do Componente de Conhecimento Específico é negativo ( – 0,48 ).
Nota -se pelo histograma ( Gráfico 3.15 ) que esta é uma distribuição com um peso maior do lado esquerdo da distribuição.
Já na região Nordeste, o coeficiente de assimetria é positivo, igual a 0,16, evidenciando que a distribuição das notas dos estudantes dessa região, no Componente de Conhecimento Específico, tem cauda mais pesada à direita.
Na sequência, os resultados verificados para cada uma das questões discursivas de Formação Geral serão apresentados, estabelecendo -se relações com os conteúdos abordados em cada uma delas.
Os comentários da Banca de docentes corretores a respeito do observado na correção das respostas dos estudantes, suas impressões e conclusões serão apresentados junto à análise de cada questão.
Cumpre esclarecer que, tendo em vista que as questões discursivas de Formação Geral são padronizadas, ou seja, constam de todas as provas, os comentários da Banca são os mesmos para todas as carreiras acadêmicas, sendo direcionados a todos os estudantes que participaram do ENADE/2011.
A seguir, serão analisados os desempenhos da Área de Engenharia – Grupo VIII nas duas questões discursivas de Formação Geral do ENADE/2011, comparando os resultados obtidos com comentários para cada questão. 3.3.1.1 Análise da Questão Discursiva 1 do Componente de Formação Geral Os dados de Engenharia – Grupo VIII, obtidos a partir das respostas à questão 1, encontram -se na Tabela 3.9 e no Gráfico 3.16.
Nessa questão – de melhor desempenho dentre as duas de Formação Geral – os alunos de todo Brasil tiveram média, 56,9.
A maior média para a questão 1 foi obtida na região Centro-Oeste ( 65,6 ), e a menor, na região Nordeste ( 41,3 ).
Quanto à variabilidade das notas, o desvio padrão de todo o Brasil foi 36,5.
O menor desvio padrão foi obtido na região Sudeste ( 33,2 ) e o maior desvio padrão foi obtido na região Nordeste ( 39,6 ).
As notas máximas e notas mínimas da questão discursiva 1 foram as mesmas para todas as regiões do Brasil, respectivamente 100,0 e 0,0.
A mediana do Brasil foi 60,0 que se repete na região Norte e Sul.
As demais medianas foram : 70,0 para as regiões Sudeste e Centro-Oeste e 30 para região Nordeste.
O Gráfico 3.16 mostra a distribuição das notas na questão discursiva 1 do Componente de Formação Geral.
Observa -se que a maior frequência corresponde aos alunos que tiveram notas compreendidas no último intervalo, ou seja, o mais alto ( 90 ; 100 ].
O coeficiente de assimetria da distribuição das notas do Componente de Conhecimento Específico é negativo ( – 0,29 ). 3.3.1.2 Comentários sobre a correção das respostas à Questão Discursiva 1 De maneira geral, a aplicação da chave de correção da questão 1 de Formação Geral não apresentou qualquer dificuldade digna de menção.
Para isso contribuíram, indubitavelmente, os ajustes feitos após a correção da amostra e a reunião entre todos os membros da Banca de docentes corretores.
As poucas dúvidas, todas pontuais, apresentadas pelos corretores, foram acompanhadas e respondidas pela coordenação e subcoordenação da correção das questões de Formação Geral, por meio da ferramenta de Gerenciamento de Dúvidas do Sistema de Correção On-line.
Não houve registro de qualquer ocorrência que pusesse em xeque o padrão de resposta ou a efetividade e a adequação da chave de correção.
Explica -se : trata -se de questão com comando claro, direto e objetivo ( solicitava -se, basicamente, três vantagens justificadas de cursos a distância ), cujas respostas foram corrigidas por meio da aplicação de um chave de correção testada e aprovada previamente.
Havia absoluta clareza quanto aos critérios de avaliação da correspondência entre as respostas dos estudantes e as possibilidades de vantagens de cursos a distância admitidas como corretas no padrão de resposta oficial, além de gradações explícitas ( e fáceis de aplicar ) dos diferentes níveis de pontuação previstos.
Felizmente, portanto, não há reparo a registrar em relação à facilidade de aplicação do padrão de resposta e da chave de correção, e nem em relação à atribuição dos diferentes níveis de pontuação previstos.
Todas as dificuldades que poderiam ter obstado a correta aplicação do padrão de resposta oficial e da respectiva chave de correção foram evitadas por meio dos ajustes feitos após a correção da amostra e fartamente debatidos com toda a Banca.
Digno de nota é que quantidade tão significativa de profissionais envolvidos na mesma tarefa – tanto para a questão 1 quanto para a 2 – tenha apresentado tão poucas dificuldades na execução da correção, em termos proporcionais.
Em suma, a correção da questão 1 da prova de Formação Geral do ENADE/2011 foi exemplarmente bem planejada, servindo -se de padrão de resposta muito bem adequado à questão proposta.
Quanto ao tema desta questão, em particular, a Banca verificou que uma parcela significativa de estudantes evocou experiências bastante concretas e próximas de sua realidade.
Houve várias respostas que indicavam uma vivência pessoal de ensino superior na modalidade Educação a Distância ( EaD ), evidenciando o tom de depoentes nos textos apresentados.
Os estudantes, em sua grande maioria, utilizaram parte considerável do espaço de 15 linhas disponíveis para a resposta – e outra parcela menos significativa dos que não o fizeram demonstraram notável capacidade de atender ao comando da questão de maneira objetiva, curta e, via de regra, correta.
Registre -se, ainda, que foram relativamente poucos os casos de respostas que tenham passado ao largo do tema em pauta na questão.
As capacidades de leitura, de compreensão do comando proposto e de expressão escrita que os estudantes avaliados na edição 2011 do ENADE foram satisfatórias.
Não obstante, seria leviano perder de vista que a qualidade dos textos redigidos em resposta às questões discursivas do Exame ainda está muito aquém do que se espera de concluintes de cursos de ensino superior de todas as regiões do país.
Quanto ao conteúdo das respostas, a Banca constatou boa capacidade, por parte da maioria dos estudantes, de compreensão do tema e do comando da questão.
Foram relativamente poucos os casos de respostas que deixaram de enumerar vantagens da modalidade EaD, e proporcionalmente escassos os estudantes que citaram vantagens não previstas no padrão de resposta.
Foi frequente, a tentativa direta de atender ao comando da questão.
Os erros mais comuns, em relação ao padrão de respostas e à grade de correção, foram fruto do desdobramento em vários “ itens ” daquilo que, de acordo com o padrão de respostas oficial, representava uma única vantagem.
Destacaram -se, neste caso, as respostas que apontavam a flexibilidade de horário e/ou local como duas vantagens distintas daquela modalidade de ensino.
Quanto aos diferentes níveis de pontuação previstos, a maior causa de baixas pontuações foi a ausência de justificativas, e mesmo de argumentação, para uma ou mais das vantagens enumeradas.
Isso demonstra não apenas a objetividade das respostas, coerente com a objetividade do comando da questão ( “ enumere três vantagens de um curso a distância ” ), mas também certa dificuldade de formulação plena de um texto, ou ao menos de parágrafos, em formato dissertativo – mesmo diante de uma média de 5 linhas disponíveis para cada vantagem a enumerar.
Os acertos mais comuns, ou seja, os “ itens ” do padrão de respostas mais frequentemente mencionados foram : ( 1 ) a flexibilidade de horário e/ou local ; ( 2 ) a capilaridade do ensino a distância ; ( 3 ) a democratização do acesso à educação de qualidade ; e ( 4 ) os custos menores que os de cursos presenciais.
Dentre as vantagens previstas no padrão de respostas que foram menos citadas, destacam -se a inclusão de pessoas com comprometimento motor, a qualificação de professores e a troca de experiências entre os participantes.
Foram muito comuns, no entanto, as menções à supostamente maior facilidade de acesso a professores ou/e tutores em cursos superiores a distância. 3.3.1.3 Análise da Questão Discursiva 2 do Componente de Formação Geral A Tabela 3.10 mostra que o desempenho dos estudantes na questão 2 ( média 46,7 ) foi inferior ao obtido na questão de número 1 ( média 56,9 ).
A região Centro-Oeste foi aquela onde a média, nessa questão, foi maior ( 55,9 ), e a de menor média foi a região Nordeste ( 31,6 ).
Quanto à variabilidade das notas, o desvio padrão de todo o Brasil foi 32,3, inferior ao obtido na questão de número 1 ( 36,5 ).
O maior desvio nessa questão foi obtido na região Nordeste ( 33,2 ), enquanto o menor foi obtido na região Centro-Oeste ( 29,5 ).
A mediana para o Brasil na questão 1 é igual a obtida na questão 2.
As notas máximas e notas mínimas da questão discursiva 2 foram as mesmas para todas as regiões do Brasil, respectivamente 100,0 e 0,0.
O Gráfico 3.17 mostra a distribuição das notas na questão discursiva 2 do Componente de Formação Geral.
Este gráfico não apresenta apenas uma moda, mas possui vários máximos locais.
O intervalo modal, o de maior frequência, é aquele que representa as questões deixadas em branco.
Desconsiderando as questões em branco, a moda da distribuição situa -se no intervalo ( 50 ; 60 ] e com ainda alta concentração no intervalo seguinte, o ( 60 ; 70 ].
O coeficiente de assimetria é -0,32.
Com exceção da região Nordeste os coeficientes de assimetria para as outras regiões são da mesma ordem de \textbf{grandeza}, sugerindo histogramas de forma semelhantes. 3.3.1.4 Comentários sobre a correção das respostas à Questão Discursiva 2 Diferentemente da questão 1, a questão 2, cujo tema aborda políticas públicas para a erradicação do analfabetismo e para o aumento da empregabilidade, permitiu mais amplitude na elaboração das respostas dos estudantes, uma vez que solicita a apresentação de uma proposta.
À falta de objetividade técnica na elaboração da questão – sobretudo no seu comando – correspondeu uma miríade de possibilidades interpretativas.
Registre -se que não raras foram as respostas que se utilizaram de lugares comuns e exortações religiosas/humanitárias/cívicas, contudo, na grande maioria dos casos, essas exortações foram usadas como um encerramento do texto e não comprometeram a resposta.
Muito frequente também foi a confusão entre política pública e política partidária, bem como alguns poucos, confundiam programas educacionais com programas educativos nas TVs.
Assim, destacaram -se como propostas/programas mais recorrentes :  prosseguimento das políticas já existentes, mas com o aumento dos investimentos, normalmente sugerindo a ampliação da bolsa família, da Educação de Jovens e Adultos ( EJA ) e da EaD ;  valorização do magistério, construção de escolas e melhoria das atuais ;  parceria do governo com empresas para manutenção de salas de aula visando a alfabetização de seus funcionários, oferecimento de estágios e redução da carga horária em troca de \textbf{isenção} de impostos ;  parcerias com igrejas e ONGs para criação de espaços de alfabetização ;  escolas itinerantes e alfabetizadores em domicílio, principalmente para pessoas com dificuldade de locomoção, como os idosos e deficientes físicos ;  erradicação do trabalho infantil ;  vinculação da bolsa família não apenas à frequência, mas também e, principalmente, ao resultado obtido pelo aluno na escola ;  revisão das políticas atuais, sendo a mais frequente a extinção da aprovação automática ;  críticas consistentes em relação a modelos didáticos considerados inadequados e desestimulantes para a educação não só de adultos, mas de pessoas de todas as idades.
Destacam -se, ainda, outras sugestões apresentadas :  \textbf{Creches} nas escolas onde os pais estão sendo alfabetizados.  Diminuição de duas horas na jornada de trabalho em empresas para funcionários não alfabetizados, para que possam frequentar a escola.  Campanhas educativas vinculadas aos meios de comunicação.  Aumento do número de escolas noturnas.  Formação específica para professores alfabetizadores.
Quanto à relação entre o analfabetismo e a empregabilidade, deve -se sublinhar que nem todos os estudantes estabeleceram claramente o vínculo entre essas duas situações sociais.
Alguns falaram separadamente de uma e de outra.
Mas a maioria fez referência à necessidade de estudo para “ conseguir um bom emprego com um bom salário ”.
Alguns, em menor número, estabeleceram de forma bastante interessante a questão histórica para a situação do Nordeste ; e, também, a relação entre a pessoa analfabeta/com pouca instrução e a desempregada/empregada em serviços mais pesados e pior remunerados, que não teriam condições de educar seus filhos que, por sua vez, também não teriam melhores oportunidades no mercado de trabalho, identificando a formação de um círculo \textbf{vicioso} e a necessidade de sua interrupção.
Muitos afirmaram que o analfabetismo não é o único responsável pelo desemprego e sim a má distribuição de renda.
Grande também foi o número de alunos que criticaram o resultado da pesquisa, afirmando que a mesma é enganosa, uma vez que considera alfabetizada a pessoa que “ desenha seu nome ”, sendo comum associarem a atual política de alfabetização com ganhos eleitorais.
Foi comum a resposta incompleta, em que o estudante fez apenas a análise das desigualdades/crítica do quadro apresentado, ou só apresentou proposta.
Alguns se limitaram a responsabilizar o governo referindo -se de forma bastante genérica à questão da “ educação ” e “ profissional ”.
Entre os equívocos que mais se repetiram, destacou -se a simples análise dos dados apresentados na tabela que consta do enunciado.
Alguns textos, inclusive, estavam corretos, sem que, no entanto, fosse respondida a pergunta.
Da mesma forma, quando os estudantes partiram para a segunda parte da questão, fizeram referência a vários programas já existentes ou simplesmente disseram que é muito importante que existam projetos para a educação.
Outro equívoco recorrente foi a análise da educação no Brasil como um todo.
Essa análise, apesar de correta, não se referia especificamente à questão do analfabetismo, o que, por conseguinte, levava a sugestões que não eram direcionadas à erradicação do analfabetismo, e sim à melhora da educação no Brasil.
Sendo assim, foram apresentadas propostas como o aumento do número de faculdades, o reforço de alunos do Ensino Médio, a criação de escolas técnicas, etc.
Por outro lado, houve um grande número de redações bem escritas e precisas, no que se refere ao que foi exigido pela questão.
Foram análises equilibradas e sensatas, correlacionando a problemática do analfabetismo com o desemprego e a apresentação de sugestões bastante consistentes. 3.3.2 Componente de Conhecimento Específico Na parte da prova relativa às questões discursivas no Componente de Conhecimento Específico ( Tabela 3.11 ), observa -se que a média foi mais baixa do que para as questões discursivas do Componente de Formação Geral.
Enquanto no Componente de Formação Geral a média para estudantes de Engenharia – Grupo VIII de todo o Brasil foi 51,9, na parte de Conhecimento Específico a média foi 16,2.
A maior média deste componente foi obtida Assim, destacaram -se como propostas/programas mais recorrentes : a continuidade das políticas já existentes, mas com o aumento dos investimentos, normalmente sugerindo a ampliação da bolsa família, da Educação de Jovens e Adultos ( EJA ) e da EaD ; valorização do magistério, construção de escolas e melhoria das atuais ; parceria do governo com empresas para manutenção de salas de aula visando a alfabetização de seus funcionários, oferecimento de estágios e redução da carga horária em troca de \textbf{isenção} de impostos ; parcerias com igrejas e ONGs para criação de espaços de alfabetização ; escolas itinerantes e alfabetizadores em domicílio, principalmente para pessoas com dificuldade de locomoção, como os idosos e deficientes físicos ; erradicação do trabalho infantil ; vinculação da bolsa família não apenas à frequência, mas também e, principalmente, ao resultado obtido pelo aluno na escola ; revisão das políticas atuais, sendo a mais frequente a extinção da aprovação automática ; críticas consistentes em relação a modelos didáticos considerados inadequados e desestimulantes para a educação não só de adultos, mas de pessoas de todas as idades.
Destacam -se, ainda, outras sugestões apresentadas : \textbf{Creches} nas escolas onde os pais estão sendo alfabetizados.
Diminuição de duas horas na jornada de trabalho em empresas para funcionários não alfabetizados, para que possam frequentar a escola.
Campanhas educativas vinculadas aos meios de comunicação.
Aumento do número de escolas noturnas.
Formação específica para professores alfabetizadores.
Quanto à relação entre o analfabetismo e a empregabilidade, deve -se sublinhar que nem todos os estudantes estabeleceram claramente o vínculo entre essas duas situações sociais.
Alguns falaram separadamente de uma e de outra.
Mas a maioria fez referência à necessidade de estudo para “ conseguir um bom emprego com um bom salário ”.
Alguns, em menor número, estabeleceram de forma bastante interessante a questão histórica para a situação do Nordeste ; e, também, a relação entre a pessoa analfabeta/com pouca instrução e a desempregada/empregada em serviços mais pesados e pior remunerados, que não teriam condições de educar seus filhos que, por sua vez, também não teriam melhores oportunidades no mercado de trabalho, identificando a formação de um círculo \textbf{vicioso} e a necessidade de sua interrupção.
Muitos afirmaram que o analfabetismo não é o único responsável pelo desemprego e sim a má distribuição de renda.
Grande também foi o número de alunos que criticaram o resultado da pesquisa, afirmando que a mesma é enganosa, uma vez que considera alfabetizada a pessoa que “ desenha seu nome ”, sendo comum associarem a atual política de alfabetização com ganhos eleitorais.
Foi comum a resposta incompleta, em que o estudante fez apenas a análise das desigualdades/crítica do quadro apresentado, ou só apresentou proposta.
Alguns se limitaram a responsabilizar o governo referindo -se de forma bastante genérica à questão da “ educação ” e “ profissional ”.
Entre os equívocos que mais se repetiram, destacou -se a simples análise dos dados apresentados na tabela que consta do enunciado.
Alguns textos, inclusive, estavam corretos, sem que, no entanto, fosse respondida a pergunta.
Da mesma forma, quando os estudantes partiram para a segunda parte da questão, fizeram referência a vários programas já existentes ou simplesmente disseram que é muito importante que existam projetos para a educação.
Outro equívoco recorrente foi a análise da educação no Brasil como um todo.
Essa análise, apesar de correta, não se referia especificamente à questão do analfabetismo, o que, por conseguinte, levava a sugestões que não eram direcionadas à erradicação do analfabetismo, e sim à melhora da educação no Brasil.
Sendo assim, foram apresentadas propostas como o aumento do número de faculdades, o reforço de alunos do Ensino Médio, a criação de escolas técnicas, etc.
Por outro lado, houve um grande número de redações bem escritas e precisas, no que se refere ao que foi exigido pela questão.
Foram análises equilibradas e sensatas, correlacionando a problemática do analfabetismo com o desemprego e a apresentação de sugestões bastante consistentes. 3.3.2 Componente de Conhecimento Específico Na parte da prova relativa às questões discursivas no Componente de Conhecimento Específico ( Tabela 3.11 ), observa -se que a média foi mais baixa do que para as questões discursivas do Componente de Formação Geral.
Enquanto no Componente de Formação Geral a média para estudantes de Engenharia – Grupo VIII de todo o Brasil foi 51,9, na parte de Conhecimento Específico a média foi 16,2.
A maior média deste componente foi obtida pelos estudantes da região Sudeste ( 23,9 ), e a menor, pelos da região Nordeste ( 6,6 ).
Quanto à variabilidade das notas, o desvio padrão de todo o Brasil foi 16,0.
O maior desvio padrão foi encontrado na região Sudeste ( 17,6 ), e o menor, na região Nordeste ( 11,5 ).
A maior nota máxima foi obtida na região Centro-Oeste ( 73,3 ) enquanto a menor nota máxima foi encontrada na região Norte ( 48,3 ).
Além disso, a nota mínima ( 0,0 ) foi obtida por alunos de todas as regiões do Brasil, sem exceção.
A mediana do Brasil como um todo foi 13,3, que não se repetiu em nenhuma região.
As medianas nas regiões foram : 11,7 na região Norte, 0,0 na região Nordeste, 25,0 na região Sudeste, 16,7 na região Sul e 20,0 na região Centro-Oeste.
O Gráfico 3.18 representa a distribuição das notas nas questões discursivas no Componente de Conhecimento Específico.
Essa distribuição é unimodal com moda no intervalo de ( 0 ; 10 ].
Nota -se uma tendência decrescente a partir do intervalo [ 0 ; 10 ] até o intervalo ( 70 ; 80 ].
Não foram registradas notas maiores que 80.
O coeficiente de assimetria da distribuição das notas do Componente de Conhecimento Específico é positivo ( 0,72 ). 3.3.2.1 Análise da Questão Discursiva 3 do Componente de Conhecimento Específico Na questão 3, cujos resultados aferidos encontram -se descritos na Tabela 3.12, a média dos estudantes de todo o Brasil foi 18,3.
A menor média nessa questão foi obtida pelos alunos da região Nordeste ( 6,0 ), enquanto a maior média foi obtida na região Sudeste ( 33,3 ).
Quanto à variabilidade das notas, o desvio padrão de todo o Brasil foi 26,2.
O maior desvio padrão foi obtido na região Sudeste ( 32,4 ), enquanto o menor foi obtido na região Norte ( 15,1 ).
A nota máxima, 100,0 pontos, foi alcançada por pelo menos um aluno das regiões Sudeste, Sul e Centro-Oeste.
A região onde ocorreu a menor nota máxima foi a Norte ( 85,0 ).
A mediana do Brasil como um todo foi 0,0 e foi a mesma nas regiões Norte e Nordeste.
Nas demais regiões a mediana foi : 25,0 na região Sudeste, 5,0 na região Sul e 10,0 na região Centro-Oeste.
A nota mínima ( 0,0 ) foi obtida em todas as regiões do Brasil, sem exceção.
O Gráfico 3.19 mostra a distribuição das notas na questão discursiva 3, do Componente de Conhecimento Específico.
Essa distribuição é unimodal com moda nas questões em branco.
Desconsiderando as questões em branco, as notas se concentram na nota zero e no intervalo ( 0 ; 10 ].
É uma distribuição altamente assimétrica com coeficiente 1,34. 3.3.2.2 Comentários sobre a correção das respostas à Questão Discursiva 3 Quanto à questão 3, que trata das formas de implantação de floresta ( talhadia, autofuste e sistemas agroflorestais ), ocorreram respostas variadas.
Neste assunto houve um desempenho superior quando comparado às respostas da questão 4, por se tratar de matéria essencial à formação do Engenheiro Florestal.
Quanto às notas zero e brancas, houve um bom número delas.
O que normalmente caracterizou a nota zero foi a redação de textos evasivos, que repetiram o enunciado da questão e que falaram sobre a importância do eucalipto na silvicultura brasileira.
Não tocaram no assunto da talhadia, auto fuste e nem mesmo dos sistemas agroflorestais, como no exemplo a seguir : “ (... ) o eucaplito é muito importante na silvicultura brasileira,(... ) ” Foi percebido que muitos confundiram talhadia e autofuste, trocando conceitos, mas souberam descrever os sistemas e suas características, ou seja, apenas confundiram os nomes.
Outro grupo, além de trocar os nomes, confundiu e misturou as características dos dois sistemas.
Os sistemas mais destacados na silvicultura brasileira são a talhadia e o auto fuste.
Um bom percentual de estudantes destacou que o sistema agroflorestal é o mais importante e difundido, e souberam caracterizar este sistema, melhor que os dois demais.
Um estudante afirmou que “ este é o sistema que deveria ser sempre preferido ”.
Houve, também, confusão com as técnicas de manejo que podem ser aplicadas com os sistemas de implantação.
Muitos falaram em desbaste, em desrama, combate à formiga e até em exploração florestal, fora do contexto dos sistemas de implantação.
Outros, ainda, descreveram os passos para se implantar uma floresta.
Quanto à grafia, houve erros, mas nada a se destacar. 3.3.2.3 Análise da Questão Discursiva 4 do Componente de Conhecimento Específico A Tabela 3.13 contém as informações relativas à questão 4 do conjunto de questões do Componente de Conhecimento Específico.
O desempenho dos estudantes de todo o Brasil nesta questão foi inferior ao desempenho na questão de número 3 e na de número 5.
A média geral do Brasil foi 5,5, sendo a menor média registrada na região Nordeste ( 1,9 ) e a maior na região Sudeste ( 8,7 ).
A nota mínima ( 0,0 ) e a mediana ( 0,0 ) foram atingidas em todas as regiões, sem exceção.
A nota máxima ( 80,0 ) do Brasil se repete na região Centro-Oeste.
A nota máxima das demais regiões foram : 50,0 na região Norte, 55,0 na região Nordeste e 65,0 nas regiões Sudeste e Sul.
O Gráfico 3.20, representa a distribuição de notas na questão discursiva 4, no Componente de Conhecimento Específico.
A nota zero e o intervalo ( 0 ; 10 ] foram o de igual e maior frequência no gráfico ( cerca de 20 \% ), não se considerando a categoria “ em branco ” ( mais de 40 \% ).
Não houve notas acima de 80,0 nesta questão.
Observa -se, ainda, que uma maior proporção de estudantes deixou essa questão em branco, comparando -se à questão anterior.
O coeficiente de assimetria é 2,7.
Os coeficientes para as regiões são todos positivos e da mesma ordem de \textbf{grandeza}, sugerindo histogramas semelhantes com alta concentração em provas em branco ( nota zero ). 3.3.2.4 Comentários sobre a correção das respostas à Questão Discursiva 4 Parte a ) – Que trata da qualidade da matéria prima para compactação de resíduos agroflorestais.
Foi percebido um baixo a moderado índice de acerto integral das respostas.
Foi perceptível que muitos dos que fizeram a prova não entenderam o item a, ou seja, o que exatamente a questão solicitava.
Esta questão pode ser considerada de moderado grau de dificuldade, devido às inúmeras possibilidades de combinações de respostas.
No gabarito final elaborado pela Comissão Assessora ENADE para a Área, a resposta para este item, estava mais associada às características que deveriam ter a matéria prima para o processo de compactação destes resíduos para fins energéticos.
Ficou evidente, também, que muitos estudantes não tinham noções a respeito deste tema, que trata de compactação de resíduos para fins energéticos, pelo tipo de respostas dadas.
Com isso, houve muitas divagações.
Poucos citaram as palavras “ peletização ” e “ briquetagem ”, apesar do texto deixar claro que a compactação ali descrita, tinha por objetivo os fins energéticos.
Neste item, o percentual baixo de acerto integral verificado pode também ser associado ao fato de se tratar de um conhecimento técnico especifico que muitos revelaram não ter.
Um grande número de estudantes manifestou por escrito que não teve esta informação na sua formação, fato que deve ser observado pelo MEC.
Parte b ) – Que trata das vantagens da compactação de resíduos.
Pelo tipo de gabarito dado, a resposta deveria estar mais associada à concentração de energia por unidade volumétrica e suas vantagens.
Nas respostas, não muitos responderam segundo o gabarito estabelecido e descreveram as vantagens silviculturais e ambientais.
Muitos “ relacionaram a resposta a aspectos de conservação do solo, citando a peletização ” e a “ briquetagem ” como uma desvantagem.
Como exemplo um estudante escreveu : “ a retirada dos resíduos de um povoamento florestal, para ser compactados, deixariam as áreas limpas, com menor indecência de incêndios, pragas, doenças e retirando nutriente que seriam, pela decomposição, incorporados ao solo ”.
Alguns estudantes confundiram compactação de solo com compactação de resíduos.
Parte c ) – Que trata da compactação dos resíduos agroflorestais e o meio ambiente.
Este foi o item no qual houve melhor desempenho dos estudantes, uma vez que é algo mais fácil de visualizar, além de ser um tema atual em nossa sociedade.
Os estudantes demonstraram entendimento sobre as consequências da exposição desta matéria orgânica em decomposição em contato com o ar, com o solo e com as águas.
Esta resposta não depende tanto de um conhecimento técnico especifico.
As respostas dos estudantes tiveram um moderado a \textbf{elevado} índice de acerto. 3.3.2.5 Análise da Questão Discursiva 5 do Componente de Conhecimento Específico A Tabela 3.14 contém as informações relativas à questão 5 do conjunto do Componente de Conhecimento Específico.
O desempenho dos estudantes nessa questão foi superior ao das questões 3 e 4.
A nota média dos estudantes de todo o Brasil foi 24,8.
A maior média foi registrada na região Centro-Oeste ( 31,5 ), enquanto a menor média foi registrada na região Nordeste ( 11,9 ).
Quanto à variabilidade das notas, o desvio padrão dos alunos do Brasil, como um todo, foi 24,5.
Enquanto o maior desvio foi encontrado na região Norte ( 26,3 ), o menor foi encontrado na região Nordeste ( 20,0 ).
A nota máxima ( 100,0 ) foi alcançada apenas na região Norte, enquanto a menor nota máxima foi registrada nas regiões Nordeste e Sul ( 80,0 ).
A nota mínima foi zero para todas as regiões do Brasil.
A mediana do Brasil foi 20,0 e a região Nordeste apresentou mediana zero, ou seja, nesta região pelo menos metade dos alunos tirou nota zero na questão.
O Gráfico 3.21 representa a distribuição das notas da questão discursiva 5 no Componente de Conhecimento Específico.
Destaca -se o grande número de estudantes que deixaram a questão 5 em branco ou obtiveram nota zero tendo respondido a questão.
O coeficiente de assimetria foi positivo, 0,42. 3.3.2.6 Comentários sobre a correção das respostas à Questão Discursiva 5 Parte a e b ) – Importância do Manejo Florestal Comunitário e o Potencial Multiplicador da exploração comunitária.
Os estudantes tiveram grande dificuldade de responder esta questão em ambos os itens, visto que muitos não souberam nem diferenciar o manejo florestal comunitário do manejo florestal convencional, ou preferiram deixar a questão em branco.
O baixo desempenho pode ser atribuído à aplicação restrita do manejo florestal comunitário, ocorrido, na maioria das vezes, apenas a região Amazônica, o que torna este método pouco abordado nas universidades do sul e sudeste, onde se encontra a maioria dos cursos de Engenharia Florestal do país.
Além disto, a viabilidade econômica, social e ambiental do manejo florestal comunitário é um tema ainda polêmico.
Como relatou um estudante : “ o manejo de madeiras na floresta tem potencial para gerar emprego ” O estudante, contudo, não expressou como obter este ganho.
A solicitação para redigir um texto dissertativo, considerando dois pontos em especial, também pode ter complicado um pouco a resposta dos estudantes, principalmente pelo pequeno espaço disponibilizado para as respostas e pela necessidade de abordar com ênfase a viabilidade econômica, social e ambiental na questão.
Outro ponto que foi observado é que muitos estudantes apenas transcreveram trechos do enunciado da questão, que era muito extenso e trazia várias informações que poderiam ser usadas como resposta.
Entretanto, em algumas respostas isso até fazia sentido, devido ao tipo de “ caput ” utilizado na elaboração da questão.
É sabida a enorme importância da floresta Amazônica e da adoção de práticas sustentáveis para conservação deste bioma.
No entanto, assuntos muito específicos de determinada região e temas polêmicos não deveriam ser abordados, pois não refletem bem a formação do engenheiro florestal e sim, a ênfase dada em cada universidade em função da demanda regional e cultural da instituição de ensino.
Assim, deveria abordar questões mais consolidadas e comuns aos cursos de Engenharia Florestal, com base na silvicultura, tecnologia da madeira ou ambiental.
As notas das respostas das duas perguntas tiveram um baixo a moderado índice de acerto. 3.3.3 Considerações Finais A Banca corretora não encontrou dificuldades em corrigir as questões 3 e 4.
Entretanto, a questão 5, por explorar um conteúdo que, em certo sentido, pode ser considerado um “ juízo de valores ”, foi mais difícil uma avaliação criteriosa, que pudesse categorizar precisamente o conteúdo apresentado pelo estudante.
Para facilitar a correção dessa questão, foi discutido entre os membros da Banca o mais adequado gabarito para correção da questão 5, embora procurando seguir as orientações dos padrões de respostas.
Mesmo com as observações acima emitidas pela Banca corretora, as perguntas da prova, acreditamos, poderão auxiliar o MEC-ENADE no processo de avaliação das Instituições de Ensino Brasileiras.

% matched lemmas: creche, elevado, grandeza, isenção, vicioso
\end{textsample}
